VIV-32 is a 32-bit load/store architecture designed for experimentation,
education, and implementation across multiple execution targets. The
architecture is intended to support a reference emulator, a QEMU target, and
a SystemVerilog or Verilog hardware implementation.

The initial version of the architecture prioritises regular instruction
decoding, straightforward toolchain development, and simple hardware
implementation over instruction density or microarchitectural sophistication.

\section{Design Goals}

\begin{itemize}
  \item Provide a clean target for assemblers, compilers, interpreters, and firmware.
  \item Keep the base ISA simple enough to implement in both software and hardware.
  \item Support a flat 32-bit physical address space.
  \item Support memory-mapped I/O.
  \item Support timer interrupts for a future supervising operating environment.
  \item Provide enough structure for an ABI, assembler, linker, and emulator.
  \item Preserve room for future extensions such as floating point, atomics, SIMD, vector, matrix, or co-processor features.
\end{itemize}

\section{Non-Goals for VIV-32 v0.1}

\begin{itemize}
  \item Virtual memory.
  \item Paging.
  \item Segmentation.
  \item Cache architecture.
  \item Multiprocessing memory coherency.
  \item Dynamic linking.
  \item User/supervisor privilege separation.
  \item Binary compatibility with an existing architecture.
\end{itemize}

\section{Specification Status}

\begin{itemize}
  \item This document describes the architectural behaviour of VIV-32.
  \item The ABI is specified separately in the VIV-32 ABI Specification.
  \item Platform devices and concrete memory maps are specified separately in the VIV-32 Platform Specification.
  \item Toolchain object formats, assembler syntax, and linker behaviour are specified separately in the VIV-32 Toolchain Specification.
  \item VIV-32 v0.1 is provisional and may change as the emulator, assembler, and hardware implementation are developed.
\end{itemize}
